\Article{50}
Es bezeichne $X$ eine Funktion der Unbestimmten $x$ von der folgenden Form:
\[
Ax^n+Bx^{n-1}+Cx^{n-2}+\text{usw.}+Mx+N;
\]

\SourceLeaf{131}%
wobei $n$ eine positive reelle ganze Zahl, $A,B,C$ usw. reelle oder imaginäre ganze Zahlen und $m$ eine komplexe ganze Zahl ist. Auch hier nennen wir eine ganze Zahl eine \emph{Wurzel} der Kongruenz $X\equiv0\pmod m$, wenn sie, für $x$ eingesetzt, für $X$ einen durch den Modul $m$ teilbaren Wert ergibt. Lösungen durch nach dem Modul kongruente Wurzeln betrachten wir nicht als verschieden.

Ist der Modul eine Primzahl, so kann auch hier eine solche Kongruenz vom Grad $n$ nicht mehr als $n$ verschiedene Lösungen besitzen. Bezeichne $\alpha$ eine beliebige feste ganze Zahl (komplex), so lässt sich $X$ durch Division durch $x-\alpha$ bei unbestimmt gelassener Veränderlicher auf die Form $X=(x-\alpha)X'+h$ bringen, wobei $h$ eine feste ganze Zahl und $X'$ eine Funktion vom Grad $n-1$ mit ganzzahligen Koeffizienten ist. Ist nun $\alpha$ eine Wurzel der Kongruenz $X\equiv0\pmod m$, so ist offenbar $h$ durch $m$ teilbar; oder es gilt bei unbestimmt gelassener Veränderlicher $X\equiv(x-\alpha)X'\pmod m$.

Bezeichnet ebenso $\beta$ eine feste ganze Zahl und wird $X'$ auf die Form $(x-\beta)X''+h'$ gebracht, so ist $X''$ eine Funktion vom Grad $n-2$ mit ganzzahligen Koeffizienten. Wird aber $\beta$ als Wurzel der Kongruenz $X\equiv0$ vorausgesetzt, so muss sie auch $(\beta-\alpha)X'\equiv0$ und auch $X'\equiv0$ erfüllen, sofern die Wurzeln $\alpha,\beta$ inkongruent sind. Daraus schließen wir, dass auch $h'$ durch $m$ teilbar sein muss, oder dass bei unbestimmt gelassener Veränderlicher $X\equiv(x-\alpha)(x-\beta)X''\pmod m$ gilt.

Kommt auf ähnliche Weise eine dritte, zu den vorigen inkongruente Wurzel $\gamma$ hinzu, so gilt bei unbestimmt gelassener Veränderlicher $X\equiv(x-\alpha)(x-\beta)(x-\gamma)X'''$, wobei $X'''$ eine Funktion vom Grad $n-3$ mit ganzzahligen Koeffizienten ist. So lässt sich weiter fortfahren, und zugleich ist klar, dass der Koeffizient des höchsten Gliedes in jeder Funktion $=A$ ist, den wir als durch $m$ nicht teilbar voraussetzen dürfen; andernfalls müsste die Kongruenz $X\equiv0$ wesentlich einem niedrigeren Grad zugerechnet werden. Sind also $n$ inkongruente Wurzeln vorhanden, etwa $\alpha,\beta,\gamma\dotsfour\nu$, so gilt bei unbestimmt gelassener Veränderlicher
\[
X\equiv A(x-\alpha)(x-\beta)(x-\gamma)\dotsfour(x-\nu)\pmod m
\]
Das Einsetzen eines neuen, zu jeder der Zahlen $\alpha,\beta,\gamma\dotsfour\nu$ inkongruenten Wertes würde folglich für $X$ sicher einen durch $m$ nicht teilbaren Wert ergeben, woraus die Richtigkeit des Satzes von selbst folgt.

Dieser Beweis stimmt im Wesentlichen mit demjenigen überein, den wir in \emph{Disq. Ar.} Artikel~43 gegeben haben und dessen einzelne Schritte für komplexe Zahlen ebenso gelten wie für reelle.

\SourceLeaf{132}%
\Article{51}
Was im dritten Abschnitt der \emph{Disquisitiones Arithmeticae} über Potenzreste gegeben wurde, gilt größtenteils mit geringfügigen Änderungen auch in der Arithmetik komplexer Zahlen; ja, meist könnten sogar die Beweise der Sätze beibehalten werden. Damit jedoch nichts fehlt, werden wir die Hauptsätze mit kurzen Beweisen anführen. Dabei ist stets vorauszusetzen, dass der Modul eine Primzahl ist.

\textsc{Satz.} \emph{Bezeichnet $k$ eine durch den Modul $m$, mit der Norm $=p$, nicht teilbare ganze Zahl, so gilt $k^{p-1}\equiv1\pmod m$.}

\emph{Beweis.} Es mögen $a,b,c$ usw. ein vollständiges System inkongruenter Reste nach dem Modul $m$ bilden, wobei jedoch der durch $m$ teilbare Rest weggelassen ist, sodass die Anzahl dieser Zahlen, deren Sammlung wir mit $C$ bezeichnen, $=p-1$ ist. Ferner sei $C'$ die Sammlung der Produkte $ka,kb,kc$ usw. Nach Voraussetzung ist keines dieser Produkte durch $m$ teilbar; daher hat jedes einen kongruenten Rest in der Sammlung $C$. Es lässt sich also $ak\equiv a',bk\equiv b',ck\equiv c'$ usw. $(\operatorname{mod.}m)$ setzen, sodass die Zahlen $a',b',c'$ usw. selbst in der Sammlung $C$ vorkommen. Die Sammlung der Zahlen $a',b',c'$ usw. bezeichnen wir mit $C''$. Es seien $P,P',P''$ die Produkte aus den einzelnen Zahlen der Sammlungen $C,C',C''$ in derselben Reihenfolge, also
\[
\begin{aligned}
P&=abc\dotsfour\\
P'&=k^{p-1}abc\dotsfour=k^{p-1}P\\
P''&=a'b'c'\dotsfour
\end{aligned}
\]
Da die Zahlen der Sammlung $C''$ der Reihe nach den Zahlen der Sammlung $C'$ kongruent sind, gilt $P''\equiv P'$ oder $P''\equiv k^{p-1}P$. Da aber leicht zu sehen ist, dass je zwei Zahlen der Sammlung $C''$ inkongruent und daher alle voneinander verschieden sind, stimmen die Zahlen der Sammlung $C''$ notwendig mit denen der Sammlung $C$ vollständig überein, nur in anderer Reihenfolge; somit gilt $P''=P$. Also ist $(k^{p-1}-1)P$ durch $m$ teilbar. Da $m$ eine Primzahl ist, die keinen der einzelnen Faktoren von $P$ teilt, muss notwendig $k^{p-1}-1$ durch $m$ teilbar sein. Q.~E.~D.

\Article{52}
\textsc{Satz.} \emph{Bezeichnet $k$ wie im vorigen Artikel eine durch den Modul $m$ nicht teilbare ganze Zahl und $t$ den kleinsten Exponenten (außer $0$), für den $k^t\equiv1\pmod m$ gilt, so ist $t$ ein Teiler jedes anderen Exponenten $u$, für den $k^u\equiv1\pmod m$ gilt.}

\SourceLeaf{133}%
\emph{Beweis.} Wäre $t$ kein Teiler von $u$, so sei $gt$ das nächstgrößere Vielfache von $u$ über $u$, sodass $gt-u$ eine positive ganze Zahl kleiner als $t$ ist. Aus $k^t\equiv1,k^u\equiv1$ folgt $0\equiv k^{gt}-k^u\equiv k^u(k^{gt-u}-1)$ und daher $k^{gt-u}\equiv1$; es gibt also eine Potenz von $k$ mit einem Exponenten kleiner als $t$, die der Einheit kongruent ist, gegen die Voraussetzung.

Als Folgerung ergibt sich daraus, dass $t$ gewiss die Zahl $p-1$ teilt.

Solche Zahlen $k$, für die $t=p-1$ gilt, nennen wir auch hier \emph{primitive Wurzeln} für. den Modul $m$. Dass es solche wirklich gibt, werden wir nun zeigen.

\Article{53}
Man zerlege die Zahl $p-1$ in ihre Primfaktoren, sodass
\[
p-1=a^\alpha b^\beta c^\gamma\dotsfour
\]
gilt, wobei $a,b,c$ usw. verschiedene positive reelle Primzahlen bezeichnen. Es seien $A,B,C$ usw. durch $m$ nicht teilbare ganze Zahlen (komplex), die jeweils den Kongruenzen
\[
x^{\frac{p-1}a}\equiv1,\quad x^{\frac{p-1}b}\equiv1,\quad x^{\frac{p-1}c}\equiv1\quad\text{usw.}
\]
nach dem Modul $m$ \emph{nicht} genügen. Dass es solche gibt, ist aus dem Satz des Artikels~50 ersichtlich. Schließlich sei $h$ nach dem Modul $m$ dem Produkt
\[
A^{\frac{p-1}{a^\alpha}}B^{\frac{p-1}{b^\beta}}C^{\frac{p-1}{c^\gamma}}\dotsfour
\]
kongruent. Dann behaupte ich, dass $h$ eine primitive Wurzel ist.

\emph{Beweis.} Bezeichnet $t$ den Exponenten der niedrigsten der Einheit kongruenten Potenz $h^t$, so wäre, wenn $h$ keine primitive Wurzel wäre, $t$ ein echter Teiler von $p-1$, oder $\frac{p-1}t$ eine die Einheit übersteigende ganze Zahl. Offenbar liegen die reellen Primfaktoren dieser ganzen Zahl unter $a,b,c$ usw. Wir nehmen daher, was erlaubt ist, an, dass $\frac{p-1}t$ durch $a$ teilbar ist, und setzen $p-1=atu$. Wegen $h^t\equiv1$ gilt dann auch $h^{tu}\equiv1$, oder
\[
A^{\frac{p-1}{a^\alpha}\pdot\frac{p-1}a}B^{\frac{p-1}{b^\beta}\pdot\frac{p-1}a}C^{\frac{p-1}{c^\gamma}\pdot\frac{p-1}a}\dotsfour\equiv1
\]
Offenbar ist aber $\frac{p-1}{ab^\beta}$ eine ganze Zahl, also
\[
B^{\frac{p-1}{b^\beta}\pdot\frac{p-1}a}=(B^{p-1})^{\frac{p-1}{ab^\beta}}\equiv1
\]
und ebenso
\SourceLeaf{134}%
$C^{\frac{p-1}{c^\gamma}\pdot\frac{p-1}a}\equiv1$ und so fort; daher muss $A^{\frac{p-1}{a^\alpha}\pdot\frac{p-1}a}\equiv1$ gelten.
Nun bestimme man eine positive ganze Zahl $\lambda$ derart, dass
\[
\lambda b^\beta c^\gamma\dotsfour\equiv1\pmod a
\]
gilt. Dies ist möglich, da die Primzahl $a$ die Zahl $b^\beta c^\gamma\dotsfour$ nicht teilt; wir setzen $\lambda b^\beta c^\gamma\dotsfour=1+a\mu$. Offenbar gilt $A^{\lambda\pdot\frac{p-1}{a^\alpha}\pdot\frac{p-1}a}\equiv1$; oder, da $\lambda\pdot\frac{p-1}{a^\alpha}\pdot\frac{p-1}a=(1+a\mu)\frac{p-1}a=(p-1)\mu+\frac{p-1}a$, haben wir $A^{(p-1)\mu}\pdot A^{\frac{p-1}a}\equiv1$. Weil von selbst $A^{(p-1)\mu}\equiv1$ gilt, folgt daraus auch $A^{\frac{p-1}a}\equiv1$, gegen die Voraussetzung. Die Annahme, dass $t$ ein echter Teiler von $p-1$ ist, kann also nicht bestehen, und somit ist $h$ notwendig eine primitive Wurzel.

